% =====================================================================
%  Zusammenfassung LEAS (HS25) – LaTeX-Quelle
%  Kompilieren: pdflatex (2x) oder latexmk -pdf main.tex
%  In VS Code mit der Extension "LaTeX Workshop": main.tex öffnen,
%  speichern (Ctrl+S) baut automatisch, Ctrl+Alt+V zeigt das PDF.
% =====================================================================
\documentclass[9pt]{extarticle}

% ---------------------------------------------------------------------
% Pakete
% ---------------------------------------------------------------------
\usepackage[utf8]{inputenc}
\usepackage[T1]{fontenc}
\IfFileExists{ngerman.ldf}{\usepackage[ngerman]{babel}}{%
  \renewcommand{\contentsname}{Inhaltsverzeichnis}}
\usepackage{helvet}
\renewcommand{\familydefault}{\sfdefault}

\usepackage[a4paper,landscape,left=0.5cm,right=0.5cm,top=0.8cm,bottom=1.0cm,footskip=0.55cm]{geometry}
\usepackage{amsmath,amssymb}
\usepackage{graphicx}
\usepackage{multicol}
\usepackage{booktabs}
\usepackage{array}
\usepackage[dvipsnames]{xcolor}
\usepackage{titlesec}
\usepackage[most]{tcolorbox}
\usepackage{fancyhdr}
\usepackage{lastpage}
\usepackage[colorlinks=true,linkcolor=blue!60,urlcolor=blue!60]{hyperref}

\graphicspath{{figures/}}

% ---------------------------------------------------------------------
% Metadaten (hier anpassen)
% ---------------------------------------------------------------------
\newcommand{\doctitle}{Zusammenfassung - LEAS}
\newcommand{\docauthors}{Jan Vogel}
\newcommand{\docdate}{29. Januar 2026}
\newcommand{\docversion}{v1.0}

% ---------------------------------------------------------------------
% Layout: Spalten, Abstände
% ---------------------------------------------------------------------
\setlength{\columnsep}{12pt}
\setlength{\columnseprule}{0.4pt}
\renewcommand{\columnseprulecolor}{\color{black!40}}
\setlength{\parindent}{0pt}
\setlength{\parskip}{1pt}
\setcounter{secnumdepth}{3}
\setcounter{tocdepth}{3}

% Abstände um abgesetzte Formeln
\setlength{\jot}{7pt}
\allowdisplaybreaks  % Formelblöcke dürfen über Spalten umbrechen
\AtBeginDocument{%
  \setlength{\abovedisplayskip}{4pt}\setlength{\belowdisplayskip}{4pt}%
  \setlength{\abovedisplayshortskip}{4pt}\setlength{\belowdisplayshortskip}{4pt}}

% ---------------------------------------------------------------------
% Überschriften als farbige Balken
% ---------------------------------------------------------------------
\definecolor{secbg}{RGB}{217,217,255}   % = blue!15
% Balken: \colorbox über die ganze Spaltenbreite, Text = Nummer + Titel
\titleformat{\section}[block]{\normalfont}{}{0pt}{\sectionbar}
\titleformat{name=\section,numberless}[block]{\normalfont}{}{0pt}{\sectionbarnonum}
\titleformat{\subsection}[block]{\normalfont}{}{0pt}{\subsectionbar}
\titleformat{\subsubsection}[block]{\normalfont}{}{0pt}{\subsubsectionbar}
\newcommand{\sectionbar}[1]{\colorbox{secbg}{\parbox{\dimexpr\linewidth-2\fboxsep\relax}{\large\strut\thesection\ #1}}}
\newcommand{\sectionbarnonum}[1]{\colorbox{secbg}{\parbox{\dimexpr\linewidth-2\fboxsep\relax}{\large\strut #1}}}
\newcommand{\subsectionbar}[1]{\colorbox{secbg}{\parbox{\dimexpr\linewidth-2\fboxsep\relax}{\footnotesize\strut\thesubsection\ #1}}}
\newcommand{\subsubsectionbar}[1]{\colorbox{secbg}{\parbox{\dimexpr\linewidth-2\fboxsep\relax}{\tiny\strut\thesubsubsection\ #1}}}

% Schriftgrössen der Balken:  section = \large, subsection = \small, subsubsection = \tiny
\titlespacing*{\section}{0pt}{4pt}{1pt}
\titlespacing*{\subsection}{0pt}{4pt}{1pt}
\titlespacing*{\subsubsection}{0pt}{4pt}{1pt}

% ---------------------------------------------------------------------
% Kopf-/Fusszeile
% ---------------------------------------------------------------------
\pagestyle{fancy}
\fancyhf{}
\renewcommand{\headrulewidth}{0pt}
\renewcommand{\footrulewidth}{0.4pt}
\fancyfoot[L]{\tiny \docauthors\ --- Zusammenfassung --- LEAS --- \docversion}
\fancyfoot[R]{\tiny \thepage/ \pageref{LastPage}}

% ---------------------------------------------------------------------
% Eigene Umgebungen
% ---------------------------------------------------------------------
% Symboltabelle:  Symbol & Bezeichnung & [Einheit] \\
\setlength{\heavyrulewidth}{0.4pt}
\setlength{\lightrulewidth}{0.4pt}
\newenvironment{symtab}{%
  \par\noindent\small\renewcommand{\arraystretch}{0.92}%
  \setlength{\aboverulesep}{0pt}\setlength{\belowrulesep}{0pt}%
  \begin{tabular}{@{\hspace{0.8em}}c@{\hspace{1.6em}}c@{\hspace{1.6em}}c@{\hspace{0.8em}}}\toprule}
  {\bottomrule\end{tabular}\par}


% Rote Hinweisbox
\newtcolorbox{achtung}[1][Achtung]{enhanced,colback=red!4,colframe=red!65,
  coltitle=white,title=#1,fonttitle=\sffamily,arc=2mm,boxrule=0.6pt,
  left=3pt,right=3pt,top=2pt,bottom=2pt,before skip=4pt,after skip=4pt}

% Bild in Spaltenbreite (Anteil von \linewidth)
\newcommand{\fig}[2][0.9]{\par\vspace{2pt}{\centering\includegraphics[width=#1\linewidth]{#2}\par}\vspace{1pt}}

% Kurzschreibweisen
\newcommand{\ul}[1]{\underline{#1}}

% =====================================================================
\begin{document}

% ---------------------------- Titelseite -----------------------------
\begin{multicols*}{3}
  {\centering
   {\Large\bfseries \doctitle\par}\vspace{4pt}
   {\normalsize \docauthors\ ---\\ \docdate\par}}
  \vspace{10pt}
  \tableofcontents
\end{multicols*}
\newpage

% ---------------------------- Inhalt ---------------------------------
\begin{multicols*}{3}
  % !TEX root = ../main.tex
% =====================================================================
\section{Grundlagen}
% =====================================================================
\subsection{Mechanische Grundgrössen}
\begin{symtab}
m & Masse & [kg] \\
s & Strecke & [m] \\
v & Geschwindigkeit & [m/s] \\
a & Beschleunigung & [m/s$^2$] \\
F & Kraft & [N] \\
B & Impuls & [kg*m/s] \\ \midrule
J & Trägheitsmoment & [kg*m$^2$] \\
$\varphi$ & Winkel & [rad] \\
$\omega$ & Winkelgeschwindigkeit & [rad/s] \\
$\omega'$ & Winkelbeschleunigung & [rad/s$^2$] \\
M & Drehmoment & [Nm] \\
D & Drall & [kg*m$^2$/s] \\ \midrule
P & Leistung & [W] \\
E & Energie & [Ws], [Nm], [J] \\
\end{symtab}

Translatorische Bewegung:
\begin{gather*}
v = \frac{ds}{dt} \\
a = \frac{dv}{dt} = \frac{d^2s}{dt^2} \\
F = m \cdot a \\
B = m \cdot v \\
P = F \cdot v \\
E = \int p(t) \cdot dt \\
E_{kin} = \frac{mv^2}{2} \\
E = \int F(s) \cdot ds
\end{gather*}

Rotatorische Bewegung:
\begin{gather*}
\omega = \frac{d\varphi}{dt} = \frac{v}{r} \\
\omega' = \frac{d\omega}{dt} = \frac{d^2\varphi}{dt^2} \\
M = J \cdot \omega' = F \cdot r \\
D = J \cdot \omega \\
E_{kin} = \frac{J\omega^2}{2} \\
E = \int M(\varphi) \cdot d\varphi
\end{gather*}

\columnbreak
% ---------------------------------------------------------------------
\subsection{Magnetismus}
\begin{symtab}
B & Magnetfeld (magnetische Induktion) & [T] \\
H & Feldstärke & [A/m] \\
N & Anz. Wicklungen & [] \\
A & Fläche & [m$^2$] \\
L & Induktivität & [H] \\
$\mu_0$ & Natürliche Permabilität & [Vs/Am] \\
$\mu_r$ & Materialabhängige Permeabilität & [] \\
$\phi$ & Magnetischer Fluss & [Vs] \\
l, l$_{fe}$ & länge, länge Eisenkern & [m] \\
F & Kraft & [N] \\
i, I & Strom & [A] \\
\end{symtab}

\subsubsection{Lorentz-Kraft}
\fig[0.25]{lorentz_kraft}
\begin{gather*}
F = i \cdot l \cdot B
\end{gather*}

\subsubsection{Magnetischer Kreis}
\fig[0.25]{magnetischer_kreis}

Durchflutungsgesetz:
\begin{gather*}
\oint H \cdot ds = H \cdot l_{fe} = \sum I = N \cdot I \\
H = \frac{N \cdot I}{l_{fe}} \\
B = \mu_0 \cdot \mu_r \cdot H \\
\phi = B \cdot A \\
\mu_0 = 4\pi 10^{-7} \\
L = \frac{N^2 \cdot A \cdot \mu_0 \cdot \mu_r}{l_{fe}}
\end{gather*}

\columnbreak
% ---------------------------------------------------------------------
\subsection{C vs. L}
\begin{symtab}
I & Strom & [A] \\
U & Spannung & [V] \\
C & Kapazität & [F] \\
L & Induktivität & [H] \\
$\omega$ & Kreisfrequenz & [rad/s] \\
E & Energie & [Ws] \\
\end{symtab}
\fig[0.9]{c_vs_l}
{\setlength{\jot}{3pt}
\begin{gather*}
E_C = \frac{C \cdot U^2}{2} \\
I_C = C \cdot \frac{du_C}{dt} \\
U_C = \frac{1}{C} \int i_C(t) dt \\
E_L = \frac{L \cdot I^2}{2} \\
I_L = \frac{1}{L} \int u_L(t) dt \\
U_L = L \cdot \frac{di_L}{dt}
\end{gather*}}

% ---------------------------------------------------------------------
\subsection{Leistung}
\begin{symtab}
I & Strom & [A] \\
U & Spannung & [V] \\
P & Leistung & [W] \\
S & Scheinleistung & [VA] \\
Q & Blindleistung & [var] \\
T & Periodendauer & [s] \\
$\lambda$ & Leistungsfaktor & [] \\
$\eta$ & Wirkungsgrad & [\%] \\
cos $\varphi$ & & [] \\
\end{symtab}

\subsubsection{einphasig}
\begin{gather*}
p(t) = u(t) \cdot i(t) \\
P = \frac{1}{T} \int_0^T p(t)\, dt = \frac{1}{T} \int_0^T u(t) \cdot i(t)\, dt \\
S = U_{rms} \cdot I_{rms} \\
P = S \cdot \cos(\varphi) \\
Q = \sqrt{S^2 - P^2} = S \cdot \sin(\varphi) \\
\lambda = \frac{P}{S} \\
\eta = \frac{P_{abgegeben}}{P_{aufgenommen}} = \frac{P_{abgegeben}}{P_{abgegeben} + P_{Verlust}}
\end{gather*}

\subsubsection{dreiphasig}
\begin{gather*}
S = 3 \cdot U_{ph\_rms} \cdot I_{ph\_rms} = \sqrt{3} \cdot U_{vk\_rms} \cdot I_{ph\_rms} \\
\ul{S} = \frac{3}{2} \cdot \ul{U_1} \cdot \ul{I_1^*} \\
P = S \cdot cos(\varphi) \\
Q = \sqrt{S^2 - P^2}
\end{gather*}

  % !TEX root = ../main.tex
% =====================================================================
\section{Gleichrichter}
% =====================================================================
\begin{symtab}
U$_{di0}$ & DC-Spannung Diode & [V] \\
U$_{di\alpha}$ & DC-Spannung Tyristor & [V] \\
U$_{d}$ & DC-Spannung nach Kommutierung & [V] \\
U$_{dx}$ & DC-Spannungsabfall an induktivität & [V] \\
U$_{N}$ & Speissespannung am Gleichrichtereingang verkettet & [V] \\
P$_{N}$ & Gleichrichterleistung & [W] \\
S$_{N}$ & Gleichrichterscheinleistung & [VA] \\
i$_{d}$ & Strom durch Last & [A] \\
i$_{HL_x}$ & Strom durch einen Halbleiter & [A] \\
i$_{k}$ & Kommutierungsstrom & [A] \\
I$_{N}$ & Eingangsstrom & [A] \\
I$_{HL}$ & Halbleiterstrom & [A] \\
U$_{HL}$ & Halbleiterspannung & [V] \\
P$_{v}HL$ & Halbleiterverlustleistung (pro HL) & [W] \\
R$_{HL}$ & Halbleiterwiderstand & [$\Omega$] \\
L$_{K}$ & Leiterinduktivität & [H] \\
$\lambda$ & Leistungsfaktor & [ ] \\
$\omega$ & Kreisfrequenz & [rad$\frac{1}{s}$] \\
f$_{aus}$ & Rippelfrequenz & [Hz] \\
f$_{ein}$ & Taktfrequenz & [Hz] \\
t & Zeit & [s] \\
ü & Überlappungswinkel & [°] \\
dx & relativer induktiver Spannungsabfall & [ ] \\
$\varphi$ & Phasenverschiebungswinkel & [°] \\
\end{symtab}

\columnbreak
% ---------------------------------------------------------------------
\subsection{Einphasig (B2)}
\fig[0.6]{b2_gleichrichter}
\begin{gather*}
i_{HLrms} = \sqrt{\frac{1}{2}} \cdot i_d \\
i_{HLavg} = \frac{1}{2} \cdot i_d \\
\lambda = \frac{P_N}{S_N} = 0.9 \cdot \cos\alpha
\end{gather*}

\subsubsection{ohne Glättung}
\begin{gather*}
U_{di0} = \frac{2\sqrt{2}}{\pi} \cdot U_N = 0.9 \cdot U_N \\
U_{di\alpha} = U_{di0} \cdot \frac{1 + \cos\alpha}{2} \\
f_{aus} = 2 \cdot f_{ein}
\end{gather*}

\subsubsection{mit Glättung}
\begin{gather*}
U_{di\alpha} = U_{di0} \cdot \cos\alpha \\
P_{NX} = U_{di0} \cdot i_d = U_{di0} \cdot i_d \cdot \cos\alpha \\
S_N = U_N \cdot I_N = \frac{\pi}{2\sqrt{2}} \cdot P_N
\end{gather*}

\columnbreak
% ---------------------------------------------------------------------
\subsection{Dreiphasig (B6)}
\fig[0.6]{b6_gleichrichter}
\begin{gather*}
U_{di0} = \frac{3\sqrt{2}}{\pi} \cdot U_N = 1.35 \cdot U_N \\
U_{di\alpha} = U_{di0} \cdot \cos\alpha \\
i_N = \sqrt{\frac{2}{3}} \cdot i_d \\
P_{di\alpha} = U_{di0} \cdot i_d \cdot \cos\alpha \\
S_N = \sqrt{3} \cdot U_{Nrms} \cdot I_{Nrms} \\
= \sqrt{3} \cdot \frac{\pi}{3\sqrt{2}} \cdot U_{di0} \cdot \sqrt{\frac{2}{3}} \cdot i_d = 1.047 \cdot U_{di0} \cdot i_d \\
\lambda = \frac{P_N}{S_N} = 0.955 \cdot \cos\alpha \\
i_{HLrms} = \sqrt{\frac{1}{3}} \cdot i_d \\
i_{HLavg} = \frac{1}{3} \cdot i_d \\
P_{vHL} = U_{HL} \cdot i_{HLavg} + i_{HLrms}^2 \cdot R_{HL}
\end{gather*}

\columnbreak
% ---------------------------------------------------------------------
\subsection{Kommutierung}
\fig[0.6]{kommutierung_b6}
\fig[0.9]{kommutierung_stromform}
\begin{gather*}
\hat{I}_K = \frac{\hat{U}_N}{2 \cdot \omega \cdot L_K} \\
I_{k(t)} = \frac{\hat{U}_N}{2 \cdot \omega \cdot L_K} \cdot (1 - \cos(\omega \cdot t)) = \hat{I}_K \cdot (1 - \cos(\omega \cdot t))
\end{gather*}
Kommutierung beginnt bei $\alpha$ und endet bei $\alpha + \text{ü}$.
\begin{gather*}
\text{ü} = \arccos\left(\cos\alpha - \frac{I_d}{\hat{I}_K}\right) - \alpha \\
U_d = U_{di0}(\cos\alpha - dx) \\
U_{dx} = U_{di0} \cdot dx \\
dx = \frac{I_d}{2 \cdot \hat{I}_K}
\end{gather*}
dx normalerweise in \% angegeben.
\begin{gather*}
\varphi = \alpha + \frac{2}{3} \cdot \text{ü} \text{ bei } \alpha \lesssim 30^\circ \\
\varphi = \alpha + \frac{1}{2} \cdot \text{ü} \text{ bei } \alpha \sim 30^\circ - 60^\circ \\
\varphi = \alpha + \frac{1}{3} \cdot \text{ü} \text{ bei } \alpha \sim 60^\circ - 90^\circ
\end{gather*}

\columnbreak
% ---------------------------------------------------------------------
\subsection{Harmonische}
{\footnotesize\setlength{\tabcolsep}{3pt}
\renewcommand{\arraystretch}{1.0}
\noindent Bei 6-pulsigen Schaltungen kommen folgende Harmonische vor:\\
\begin{tabular*}{\linewidth}{@{}l@{\extracolsep{\fill}}llllllr@{}}
$U_d$ & n*6:      & 6.     & 12.      & 18.      & 24.      & 30.      & etc. \\
$i_N$ & n*6 $\pm$ 1: & 5., 7. & 11., 13. & 17., 19. & 23.,25. & 29.,31. & etc. \\
\end{tabular*}\\[4pt]
Bei 12-pulsigen Schaltungen kommen folgende Harmonische vor:\\
\begin{tabular*}{\linewidth}{@{}l@{\extracolsep{\fill}}llllllr@{}}
$U_d$ & n*12:      & & 12.      & & 24.      & & etc. \\
$i_N$ & n*12$\pm$1: & & 11., 13. & & 23.,25. & & etc. \\
\end{tabular*}\\[4pt]
Bei 24-pulsigen Schaltungen kommen folgende Harmonische vor:\\
\begin{tabular*}{\linewidth}{@{}l@{\extracolsep{\fill}}llllllr@{}}
$U_d$ & n*24:      & & & & 24.      & & etc. \\
$i_N$ & n*24$\pm$1: & & & & 23.,25. & & etc. \\
\end{tabular*}\par}

  % !TEX root = ../main.tex
% =====================================================================
\section{Motoren}
% =====================================================================
\subsection{GM}
\subsection{Hermonische}
\fig[0.6]{gm_fremderregt}
\begin{symtab}
U$_a$ & Ankerspannung & [V] \\
U$_e$ & Erregerspannung & [V] \\
U$_i$ & Induziertespannung & [V] \\
I$_a$ & Ankerstrom & [A] \\
I$_e$ & Erregerstrom & [A] \\
R$_a$ & Ankerwiderstand & [$\Omega$] \\
R$_e$ & Erregerwiderstand & [$\Omega$] \\
L$_a$ & Ankerinduktivität & [H] \\
L$_e$ & Erregerinduktivität & [H] \\
P$_{el}$ & Leistung & [W] \\
N$_e$ & Erregerwicklungszahl & [ ] \\
M$_x$ & Drehmoment & [Nm] \\
J & Trägheitsmoment & [xxxx] \\
$\Phi$ & Erregerfluss & [Vs] \\
c & Maschinenkonstante & [ ] \\
\end{symtab}
{\setlength{\jot}{3pt}
\begin{gather*}
U_a = R_a \cdot I_a + L_a \cdot \frac{dI_a}{dt} + U_i \\
U_e = R_e \cdot I_e + L_e \cdot \frac{dI_e}{dt} \\
M_{el} = M_{Last} + M_R + J \frac{d\omega_m}{dt}
\end{gather*}}

\columnbreak
{\setlength{\jot}{3pt}
\begin{gather*}
U_i = c \cdot \Phi \cdot \omega_m \\
M_{el} = c \cdot \Phi \cdot I_a \\
\Phi = \frac{L_e}{N_e} \cdot I_e \\
P_{el} = M_{el} \cdot \omega_m + R_e \cdot I_e^2 + R_a \cdot I_a^2
\end{gather*}
stationär
\begin{gather*}
U_a = U_i + R_a \cdot I_a = c \cdot \Phi \cdot \omega_m + R_a \cdot I_a \\
n = \omega_m \cdot \frac{60}{2 \cdot \pi}
\end{gather*}}
\fig[0.6]{gm_seriegeregt}
{\setlength{\jot}{3pt}
\begin{gather*}
I_a = I_e \\
\Phi = \frac{L_e}{N_e} \cdot I_e \\
\omega_m = \frac{U_a - (R_a + R_e) \cdot I_a}{\Phi} \\
M_{el} = c \cdot \frac{L_e}{N_e} \cdot I_a^2
\end{gather*}}

% ---------------------------------------------------------------------
\subsection{Universalmotor}
\fig[0.6]{universalmotor}
{\setlength{\jot}{3pt}
\begin{gather*}
\ul{U}_e = \ul{I} \cdot R_e + j\omega L_e \cdot \ul{I} \\
\ul{U}_a = \ul{I} \cdot R_a + j\omega L_a \cdot \ul{I} \\
\ul{U}_i = c \cdot \Phi \cdot \omega_m = c \cdot \frac{L_e}{N_e} \cdot \ul{I} \cdot \omega_m \\
\ul{U} = \ul{U}_e + \ul{U}_a + \ul{U}_i = \ul{U}_i + \ul{I}(R_A + R_e) + j\omega\ul{I}(L_a + L_e) \\
M_{el} = c \cdot \Phi \cdot |\ul{I}|
\end{gather*}}

\columnbreak
% ---------------------------------------------------------------------
\subsection{ASM}
\begin{symtab}
$\omega_{D1}$ & Statordrehfeldfrequenz & [ ] \\
$\omega_{D2}$ & Rotordrehfeldfrequenz & [ ] \\
$\omega_1$ & Netzfrequenz & [1/s] \\
$\omega_2$ & Rotor ind. Frequenz & [1/s] \\
$\omega_{syn}$ & Leerlauffrequenz & [1/s] \\
$\omega_{mech}$ & mech. Drehfrequenz & [1/s] \\
p & Polpaarzahl & [ ] \\
s & Schlupf & [ ] \\
n & Drehzahl & [rpm] \\
$\eta$ & Wirkungsgrad & [\%] \\
U$_{n1}$ & Statorspannung ind. & [V] \\
U$_{n2}$ & Rotorspannung ind. & [V] \\
$\Phi_n$ & Hauptfluss & [Vs] \\
ü & Übersetzungsverhältnis & [] \\
\end{symtab}

\subsubsection{Grundlagen}
\begin{gather*}
\omega_{D1} = \frac{\omega_1}{p} \\
\omega_{mech} = 2\pi \frac{n}{60} \\
\omega_{D2} = \omega_{D1} - \omega_{mech} \\
\omega_2 = \omega_{D2} \cdot p = \omega_1 - p \cdot \omega_{mech} = s \cdot \omega_1 \\
\omega_{syn} = \omega_{D1} \\
s = \frac{n_{syn} - n}{n_{syn}} = \frac{\omega_{syn} - \omega_{mech}}{\omega_{syn}} = \frac{\omega_1 - p \cdot \omega_{mech}}{\omega_1} = \frac{\omega_2}{\omega_1}
\end{gather*}
\fig[0.45]{asm_stator_rotor}
\begin{gather*}
U_{n1} = N_1 \cdot \frac{d\Phi_n}{dt} = \omega_1 \cdot N_1 \cdot \Phi_n = \frac{N_1}{N_2} \cdot \frac{\omega_1}{\omega_2} \cdot U_{n2} = \frac{\text{ü}}{s} \cdot U_{n2} \\
U_{n2} = N_2 \cdot \frac{d\Phi_n}{dt} = \omega_2 \cdot N_2 \cdot \Phi_n \\
I_1 = \frac{1}{\text{ü} \cdot I_2} % TODO prüfen: korrekt wäre I_1 = I_2 / ü
\end{gather*}

\columnbreak
\subsubsection{Betriebgrenzwerte}
\fig[0.75]{asm_betriebsgrenzwerte}
\begin{gather*}
Stillstand: n = 0, s = 1, \omega_2 = \omega_1, U_{n2} = \frac{U_{n1}}{\text{ü}} \\
n_{syn}: n = n_{syn}, s = 0, \omega_2 = 0, U_{n2} = 0
\end{gather*}

\subsubsection{Leistung}
\begin{symtab}
P$_1$ & Statorleistung & [W] \\
P$_2$ & Rotorleistung & [W] \\
P$_\delta$ & Luftspaltleistung & [W] \\
P$_{Fe}$ & Eisenleistung & [W] \\
P$_{mech}$ & Mechanische Leistung & [W] \\
P$_s$ & Schlupfleistung & [W] \\
$\eta$ & Wirkungsgrad & [\%] \\
s & Schlupf & [ ] \\
\end{symtab}
\begin{gather*}
P_1 = 3 \cdot U_{n1} \cdot I_1 = \frac{P_2}{s} \\
P_2 = 3 \cdot U_{n2} \cdot I_2 = s \cdot P_\delta \\
P_\delta = P_1 : P_\delta = P_1 - 3 \cdot R_1 \cdot I_1^2 - P_{Fe} \text{ wenn } R_1 \& P_{Fe} \neq 0 \\
P_\delta = P_2 - P_{mech} \\
P_{mech} = P_\delta - P_s = (1 - s) \cdot P_\delta = P_1 - P_2
\end{gather*}
Wirkungsgrad
\begin{gather*}
\eta = \frac{P_{mech}}{P_1} < \frac{P_{mech}}{P_\delta} = 1 - s \quad f\ddot{u}r \quad 0 < s < 1 \\
\eta = \frac{P_1}{P_{mech}} < \frac{P_\delta}{P_{mech}} = \frac{1}{1 - s} \quad f\ddot{u}r \quad s < 0
\end{gather*}

\columnbreak
\fig[0.9]{asm_betriebsarten}

\subsubsection{Ersatzschema}
\begin{symtab}
R$_1$ & Statorwiderstand & [$\Omega$] \\
R$_2$ & Rotorwiderstand & [$\Omega$] \\
R$'_2$ & Rotorwiderstand auf Statorseite & [$\Omega$] \\
R$_s$ & Läuferwiderstand & [$\Omega$] \\
R$'_s$ & Läuferwiderstand auf Statorseite & [$\Omega$] \\
R$_{Fe}$ & Streuwiderstand & [$\Omega$] \\
L$_{\sigma 1}$ & Statorinduktivität & [H] \\
L$_{\sigma 2}$ & Rotorinduktivität & [H] \\
L$'_{\sigma 2}$ & Rotorinduktivität auf Statorseite & [H] \\
L$_h$ & Streuinduktivität & [H] \\
U$_1$ & Statorspannung & [V] \\
U$_2$ & Rotorspannung & [V] \\
U$'_2$ & Rotorspannung auf Statorseite & [V] \\
ü & Übersetzungsverhältnis & [] \\
$\psi_1$ & Statorfluss & [Vs] \\
$\psi_2$ & Rotorfluss & [Vs] \\
p & Polparzahl & [] \\
\end{symtab}
\fig[0.9]{asm_ersatzschema}
{\setlength{\jot}{1pt}
\begin{gather*}
U'_2 = \text{ü} \cdot \ul{U}_2 \\
I'_2 = \frac{I_2}{\text{ü}} \\
R'_2 = \text{ü}^2 \cdot R_2 \\
R'_s = \text{ü}^2 \cdot R_s \\
R'_s + R'_2 = \frac{R'_2}{s} \\
L'_{\sigma 2} = \text{ü}^2 \cdot L_{\sigma 2}
\end{gather*}}

\subsubsection{Drehmoment}
\begin{symtab}
M$_{el}$ & elektrisches Drehmoment & [Nm] \\
M$_k$ & Kippmoment & [Nm] \\
M$_{Reibung}$ & Reibungsmoment & [Nm] \\
P$_{mech}$ & mechanische Leistung & [W] \\
P$_\sigma$ & Luftspaltleistung & [W] \\
s$_k$ & Kippschlupf & [] \\
\end{symtab}
\fig[0.9]{asm_drehmomentkennlinie}
{\setlength{\jot}{1pt}
\begin{gather*}
M_{el} = \frac{P_{mech}}{\omega_{mech}} = \frac{(1 - s) \cdot P_\sigma}{(1 - s) \cdot \omega_{D1}} = \frac{P_\sigma \cdot p}{\omega_1} \\
M_{el} = \frac{3 \cdot p}{\omega_1} \cdot \frac{U_1^2}{(R_1 + \frac{R'_2}{s})^2 + (\omega_1 \cdot L_{\sigma 1} + \omega_1 \cdot L'_{\sigma 2})^2} \cdot \frac{R'_2}{s} \\
P_{mech} = 3 \cdot I_2^2 \cdot R_s \\
M_{el} \cong \frac{3 \cdot p}{\omega_1} \cdot \frac{U_1^2}{R'_2} \cdot s
\end{gather*}}
Folgende für Kipppunkt
{\setlength{\jot}{1pt}
\begin{gather*}
\frac{dM_{el}}{ds} = 0 \\
s_k = \frac{R'_2}{\sqrt{R_1^2 + (\omega_1 \cdot L_{\sigma 1} + \omega_1 \cdot L'_{\sigma 2})^2}} \\
\text{bei } R_1 = 0 \rightarrow s_k = \frac{R'_2}{\omega_1 \cdot (L_{\sigma 1} + L'_{\sigma 2})} \\
M_k = \frac{3 \cdot p}{2 \cdot \omega_1} \cdot \frac{U_1^2}{R_1 + \omega_1 \cdot L_{\sigma 1} + \omega_1 \cdot L'_{\sigma 2}} % TODO prüfen: im Nenner fehlt die Wurzel
\\
\text{bei } R_1 = 0 \rightarrow M_k = \frac{3 \cdot p}{2(L_{\sigma 1} + L'_{\sigma 2})} \cdot \frac{U_1^2}{\omega_1^2} \\
M = M_{el} - M_{Reibung} \\
\frac{M_{el}}{M_k} = \frac{2}{\frac{s}{s_k} + \frac{s_k}{s}}
\end{gather*}}

\columnbreak
\subsubsection{Flussverkettung rechnen}
\begin{symtab}
I$_{1n}$ & Strom Nennbetrieb & [A] \\
I$_{10}$ & Strom Leerlauf & [A] \\
P & Wirkleistung El. aufgenommen & [W] \\
\end{symtab}
\begin{achtung}
Spannungen sind Einphasig
\end{achtung}
\begin{gather*}
\Psi_x = \frac{\ul{U}_{Lx}}{j\omega_x} \\
\ul{I}_{1n} = \|\hat{I}_{1n}\| \cdot (\cos(\varphi) + j \cdot \sin(\varphi)) \\
\cos(\varphi) = \frac{P}{S} = \frac{P}{3 \cdot U_1 \cdot I_1}
\end{gather*}

\subsubsection{Zeiger Dynamisch Stator ($\alpha,\beta-KoS$)}
{\setlength{\jot}{3pt}
\begin{gather*}
\ul{u}_1 = R_1 \cdot \ul{i}_1 + \frac{d\ul{\psi}_1}{dt} \\
\ul{u}_2 = R_2 \cdot \ul{i}_2 + \frac{d\ul{\psi}_2}{dt} - j \cdot p \cdot \omega_{mech} \cdot \ul{\psi}_2 \\
\ul{\psi}_1 = L_1 \cdot \ul{i}_1 + L_h \cdot \ul{i}_2 \\
\ul{\psi}_1 = (L_h + L_{\sigma 1}) \cdot \ul{i}_1 + L_h \cdot \ul{i}_2 \\
\ul{\psi}_1 = L_h \cdot (\ul{i}_1 + \ul{i}_2) + L_{\sigma 1} \cdot \ul{i}_1 \\
\ul{\psi}_2 = L_2 \cdot \ul{i}_2 + L_h \cdot \ul{i}_1 \\
\ul{\psi}_2 = (L_h + L_{\sigma 2}) \cdot \ul{i}_2 + L_h \cdot \ul{i}_1 \\
\ul{\psi}_2 = L_h \cdot (\ul{i}_1 + \ul{i}_2) + L_{\sigma 2} \cdot \ul{i}_2 \\
M = \frac{3p}{2} \cdot Im(\ul{\psi}_1^* \cdot \ul{i}_1) = \frac{3p}{2} \frac{L_h}{L_1} \cdot Im(\ul{\psi}_2^* \cdot \ul{i}_1) = \frac{3p}{2} \frac{L_h}{L_2} \cdot Im(\ul{\psi}_1^* \cdot \ul{i}_2) \\
\omega_{mech} = \frac{1}{J} \int (M - M_{Last}) dt
\end{gather*}}

\subsubsection{Zeiger Dynamisch Rotor (I, Q-KoS)}
{\setlength{\jot}{3pt}
\begin{gather*}
\ul{u}_1 = R_1 \cdot \ul{i}_1 + \frac{d\ul{\psi}_1}{dt} + j \cdot \omega_k \cdot \ul{\psi}_1 \\
\ul{u}_2 = R_2 \cdot \ul{i}_2 + \frac{d\ul{\psi}_2}{dt} + j \cdot (\omega_k - p \cdot \omega_{mech}) \cdot \ul{\psi}_2 \\
\ul{\psi}_1 = L_1 \cdot \ul{i}_1 + L_h \cdot \ul{i}_2 \\
\ul{\psi}_2 = L_2 \cdot \ul{i}_2 + L_h \cdot \ul{i}_1 \\
M = \frac{3p}{2} \cdot Im(\ul{\psi}_1^* \cdot \ul{i}_1) = \frac{3p}{2} \frac{L_h}{L_1} \cdot Im(\ul{\psi}_2^* \cdot \ul{i}_1) = \frac{3p}{2} \frac{L_h}{L_2} \cdot Im(\ul{\psi}_1^* \cdot \ul{i}_2) \\
\omega_{mech} = \frac{1}{J} \int (M - M_{Last}) dt
\end{gather*}}

\columnbreak
% ---------------------------------------------------------------------
\subsection{SM}
\subsubsection{Allgemein}
\begin{symtab}
M$_{el}$ & Elektrisches Drehmoment & [Nm] \\
M$_{mech}$ & Mechanisches Moment & [Nm] \\
$\omega_{el}$ & elektrische Kreisfrequenz & [Hz] \\
$\omega_{mech}$ & mechanische Kreisfrequenz & [Hz] \\
p & Polpaarzahl & [] \\
U$_1$ & Klemmenspannung & [V] \\
U$_p$ & Polradspannung ind. & [V] \\
U$_{R_1}$ & Widerstandsspannung & [V] \\
U$_{L_a}$ & Leiterinduktivitätspannung & [V] \\
U$_{L_\sigma}$ & Induktivität Primärseite Spannung & [V] \\
U$_{L_d}$ & Kombinierte Leiterinduktivitätspannung & [V] \\
$\Delta\ul{U}$ & Kombinierte Leiterspannung & [V] \\
U$_{1d}$ & D-Komponente der Spannung & [V] \\
U$_{1q}$ & Q-Komponente der Spannung & [V] \\
R$_1$ & Leiterwiderstand & [$\Omega$] \\
X$_d$ & Spulenwiderstand & [$\Omega$] \\
L$_a$ & Leiterinduktivität & [H] \\
L$_\sigma$ & Induktivität Primärseite & [H] \\
L$_d$ & Kombinierte Induktivität Primärseite & [H] \\
I$_1$ & Leiterstrom & [A] \\
I$_{1d}$ & D-Komponente der Stroms & [A] \\
I$_{1q}$ & Q-Komponente der Stroms & [A] \\
$\psi_p$ & Phasenwinkel & [] \\
$\vartheta$ & Polradwinkel & [] \\
\end{symtab}
\fig[0.9]{sm_ersatzschaltung}
{\setlength{\jot}{3pt}
\begin{gather*}
M_{mech} = p \cdot M_{el} \\
\omega_{mech} = \frac{\omega_{el}}{p} \\
\ul{U}_p = j\omega_1 \cdot \ul{\psi}_p \\
\ul{U}_{L_d} = L_d \cdot \frac{di_1}{dt} \\
\ul{U}_1 = \ul{I}_1(R_1 + j\omega L_d) + \ul{U}_p = \Delta\ul{U} + \ul{U}_p \\
\Delta\ul{U} = jX_d \cdot \ul{I}_1 + R_1 \cdot \ul{I}_1 \\
X_d = \omega L_d = \omega(L_\sigma + L_a) \\
\ul{I}_1 = -j \cdot \frac{\Delta\ul{U}}{X_d}
\end{gather*}}

\columnbreak
\fig[0.9]{sm_am_netz}

\subsubsection{Drehmoment}
\fig[0.9]{sm_drehmoment}
\begin{gather*}
U_{1d} = U_1 \cdot \sin\vartheta \\
U_{1q} = U_1 \cdot \cos\vartheta \\
I_{1d} = \frac{U_1 \cdot \cos(\vartheta) - U_p}{X_d} \\
I_{1q} = \frac{U_1 \cdot \sin(\vartheta)}{X_d} \\
I_1 = \sqrt{I_{1d}^2 + I_{1q}^2}
\end{gather*}

\columnbreak
% ---------------------------------------------------------------------
\subsection{Dynamische Zeiger für SM}
\begin{symtab}
L$_d$ & Ersatzwicklung von Stator in d & [H] \\
L$_q$ & Ersatzwicklung von Stator in q & [H] \\
L$_D$ & Ersatzwicklung von Dämpferkreis in d & [H] \\
L$_Q$ & Ersatzwicklung von Dämpferkreis in q & [H] \\
L$_E$ & Ersatzwicklung von Erregerwicklung & [H] \\
L$_{qQ}$ & Kopplungsinduktivität bei $L_q$ und $L_Q$ & [H] \\
L$_{ED}$ & Kopplungsinduktivität bei $L_E$ und $L_D$ & [H] \\
L$_{dD}$ & Kopplungsinduktivität bei $L_d$ und $L_D$ & [H] \\
L$_{dE}$ & Kopplungsinduktivität bei $L_d$ und $L_E$ & [H] \\
$\Psi_d$ & Stator verketteter Fluss in d & [Vs] \\
$\Psi_q$ & Stator verketteter Fluss in q & [Vs] \\
$\Psi_D$ & Dämpferkreis verketteter Fluss in d & [Vs] \\
$\Psi_Q$ & Dämpferkreis verketteter Fluss in q & [Vs] \\
$\Psi_E$ & Erreger verketteter Fluss & [Vs] \\
$\Psi_0$ & Erreger verketteter Fluss vereinfacht & [Vs] \\
\end{symtab}
\fig[0.55]{sm_dq_wicklungen}

\subsubsection{Statorspannung und Statorfluss}
{\setlength{\jot}{2pt}
\begin{gather*}
U_{1d} = R_1 \cdot I_d + \frac{d\Psi_d}{dt} - p \cdot \omega_{mech} \cdot \Psi_q \\
U_{1q} = R_1 \cdot I_q + \frac{d\Psi_q}{dt} - p \cdot \omega_{mech} \cdot \Psi_d \\
\Psi_d = L_d \cdot I_d + L_{dE} \cdot I_E + L_{dD} \cdot I_D \\
\Psi_q = L_q \cdot I_q + L_{qQ} \cdot I_Q
\end{gather*}}

\subsubsection{Dämpferspannung und Dämpferfluss}
{\setlength{\jot}{2pt}
\begin{gather*}
0 = R_D \cdot I_D + \frac{d\Psi_D}{dt} \\
0 = R_Q \cdot I_Q + \frac{d\Psi_Q}{dt} \\
\Psi_D = L_{dD} \cdot I_d + L_{ED} \cdot I_E + L_D \cdot I_D \\
\Psi_Q = L_{qQ} \cdot I_q + L_Q \cdot I_Q
\end{gather*}}

\subsubsection{Erregerspannung und Erregerfluss}
{\setlength{\jot}{2pt}
\begin{gather*}
U_E = R_E + I_E + \frac{d\Psi_E}{dt} % TODO prüfen: gemeint ist R_E * I_E
\\
\Psi_E = L_{dE} \cdot I_d + L_E \cdot I_E + L_{ED} \cdot I_D
\end{gather*}}

\columnbreak
\subsubsection{Drehmoment}
\begin{gather*}
M = \frac{3 \cdot p}{2} \cdot (I_q \cdot \Psi_d - I_d \cdot \Psi_q) \\
\omega_{mech} = \frac{1}{J} \int (M - M_{Last}) dt
\end{gather*}

\subsubsection{Vereinfacht}
Diese Formeln gelten, falls:\\
- Die Dämpferwicklungen weggelassen werden\\
- Das Feld, verursacht durch den Erregerstrom, als konstant angenommen wird.\\
Bei Vollpolmaschienen kann weiter vereinfacht werden, da $L_d = L_q$
{\setlength{\jot}{4pt}
\begin{gather*}
\Psi_0 = L_{dE} \cdot I_E \\
U_{1d} = R_1 \cdot I_d + \frac{d\Psi_d}{dt} - p \cdot \omega_{mech} \cdot \Psi_q \\
U_{1q} = R_1 \cdot I_q + \frac{d\Psi_q}{dt} - p \cdot \omega_{mech} \cdot \Psi_d \\
\Psi_d = L_d \cdot I_d + \Psi_0 \\
\Psi_q = L_q \cdot I_q \\
I_d = \frac{\Psi_d - \Psi_0}{L_d} \\
I_q = \frac{\Psi_q}{L_q} \\
M = \frac{3 \cdot p}{2}(I_q \cdot \Psi_d - I_d \cdot \Psi_q) = \frac{3 \cdot p}{2}(I_q \cdot \Psi_0 + I_d \cdot I_q(L_d - Lq)) \\
\omega_{mech} = \frac{1}{J} \int (M - M_{Last}) dt
\end{gather*}}

% ---------------------------------------------------------------------
\subsection{Vollpol}
\begin{gather*}
P_{mech} = M_{mech} \cdot \omega_{mech} = 3(U_{1d} \cdot I_{1d} + U_{1q} \cdot I_{1q}) \\
P_{mech} = \frac{3 \cdot U_1 \cdot U_p}{X_d} \cdot \sin\vartheta = 3 \cdot U_p \cdot I_{1q} \\
\omega_{mech} = \frac{\omega_1}{p} = 2\pi f_{mech} = 2\pi \frac{n}{60} \\
M_{mech} = \frac{3 \cdot p \cdot U_1 \cdot U_p}{\omega_1 \cdot X_d} \cdot \sin\vartheta = \frac{3 \cdot p}{\omega_1} \cdot U_p \cdot I_{1p} = \frac{P_{mech}}{\omega_{mech}}
\end{gather*}

% ---------------------------------------------------------------------
\subsection{Schenkelpol}
\begin{gather*}
P_{mech} = \frac{3 \cdot U_1 \cdot U_p}{X_d} \cdot \sin(\vartheta) + \frac{3 \cdot U_1^2}{2}\left(\frac{1}{Xq} - \frac{1}{Xd}\right) \cdot \sin(s \cdot \vartheta) % TODO prüfen: üblich ist sin(2*theta)
\\
M_{mech} = M_{kipp} \cdot \sin(\vartheta)
\end{gather*}

  % !TEX root = ../main.tex
\columnbreak
% =====================================================================
\section{Transformatoren}
% =====================================================================
\begin{symtab}
U & Spannung & [V] \\
u$_1$ & Spannung Primärseite & [V] \\
u$_2$ & Spannung Sekundärseite & [V] \\
I & Strom & [A] \\
i$_1$ & strom Primärseite & [A] \\
i$_2$ & Strom Sekundärseite & [A] \\
i$_{m1}$ & magnetisierungsstrom Primärseite & [A] \\
H & Feldstärke & [A/m] \\
N & Anzahl Wicklungen & [] \\
N$_1$ & Wicklungen Primärseite & [] \\
N$_2$ & Wicklungen Sekundärseite & [] \\
$\phi$ & Magnetischer Fluss & [Vs] \\
ü & Übersetzungsverhältnis & [] \\
ü$_k$ & Spannungsfaktor damit Nennstrom fliesst & [] \\
R$_{fe}$ & Ummagnetisierungsverluste & [$\Omega$] \\
R$_1$ & Wicklungswiderstand Primärseite & [$\Omega$] \\
R$_2$ & Wicklungswiderstand Sekundärseite & [$\Omega$] \\
L$_{h1}$ & Streuinduktivität & [H] \\
L$_{\sigma 1}$ & Induktivität Primärseite & [H] \\
L$_{\sigma 1}$ & Induktivität Sekundärseite & [H] \\
\end{symtab}

\subsection{Ideal}
\fig[0.3]{trafo_ideal}
{\setlength{\jot}{3pt}
\begin{gather*}
U = N \cdot \frac{d\phi}{dt} \\
\phi = \frac{1}{N} \int U dt \\
\text{ü} = \frac{i_2}{i_1} = \frac{u_1}{u_2} = \frac{N_1}{N_2}
\end{gather*}}

\subsubsection{Magnetisierungsstrom}
\fig[0.6]{trafo_magnetisierung}
{\setlength{\jot}{3pt}
\begin{gather*}
L_{h1} = \mu_0 \cdot \mu_r \cdot \frac{A_{fe}}{l_{fe}} \cdot N_1^2 \\
i_{m1} = i_1 - \frac{N_2}{N_1} \cdot i_2
\end{gather*}}

\columnbreak
{\setlength{\jot}{3pt}
\begin{gather*}
L_{h2} = \frac{L_{h1}}{\text{ü}^2} \\
\phi = \mu_0 \cdot \mu_r \cdot A_{fe} \cdot H \\
u_1 = L_{h1} \cdot \frac{di_{m1}}{dt}
\end{gather*}}

\subsubsection{Streuung}
\fig[0.6]{trafo_streuung}

\subsection{Real}
\subsubsection{Wicklungsverluste}
\fig[0.75]{trafo_wicklungsverluste}
{\setlength{\jot}{3pt}
\begin{gather*}
u_1 = N_1 \cdot \frac{d}{dt}(\phi_h + \phi_{\sigma 1}) \\
u_2 = N_2 \cdot \frac{d}{dt}(\phi_h + \phi_{\sigma 1}) % TODO prüfen: sekundär wohl phi_sigma2
\end{gather*}}

\subsubsection{Bezug auf Primärseite}
\fig[0.9]{trafo_bezug_primaer}

\columnbreak
\fig[0.9]{trafo_t_ersatzschaltung}
{\setlength{\jot}{3pt}
\begin{gather*}
u'_2 = u_2 \text{ü} \\
i'_2 = \frac{i_2}{\text{ü}} \\
R'_2 = R_2 \text{ü}^2 \\
L'_{\sigma 2} = L_{\sigma 2} \text{ü}^2
\end{gather*}}

\subsubsection{Leerlaufversuch}
Bei den Versuchen werden komplexe Wechselstromrechnungen mit $\|Z\|$ gemacht.
\fig[0.9]{trafo_leerlauf}

\subsubsection{Kurzschlussversuch}
\fig[0.63]{trafo_kurzschluss}
ü$_k$ wird ofmals als u$_k$ in \% angegeben\par
Oft wird folgende Annahme gemacht:
\begin{gather*}
L_{\sigma 1} = L'_{\sigma 2} = \frac{L_{\sigma 1} + L'_{\sigma 2}}{2}
\end{gather*}
R$_1$ und R$_2$ lassen sich direkt messen.\par
Spetzialfall falls L$_\sigma$ nicht gegeben und L$_h$ gesucht ist.
\begin{gather*}
\omega L_h = \sqrt{Z_{Kurzschluss}^2 - R_{Gemessengesammt}^2}
\end{gather*}

\columnbreak
\subsubsection{Sättigung}
\fig[0.7]{trafo_saettigung}

  % !TEX root = ../main.tex
% =====================================================================
\section{Wechselrichter}
% =====================================================================
\begin{symtab}
U & Spannung & [V] \\
U$_{dc}$ & DC-Spannung, Zwischenkreisspannung & [V] \\
U$_a(t)$ & Wechselrichterspannung & [V] \\
U$_g$ & Gegenspannung & [V] \\
I & Strom & [A] \\
v & Harmonische (3, 5, 7, \dots) & [] \\
$\alpha$ & Zündwinkel & [°] \\
M & Modulationsgrad & [] \\
$\hat{I}_{a,v}(t)$ & Peakstrom einer Harmonischen & [A] \\
\end{symtab}

\subsection{Einphasig}
\fig[0.6]{wr_einphasig}

\columnbreak
\fig[0.6]{wr_einphasig_verlaeufe}

\subsection{Einphasige Brückenschaltung}
\fig[0.9]{wr_bruecke}

\columnbreak
\fig[0.6]{wr_bruecke_verlaeufe}

Grundfrequenztaktung:
\begin{gather*}
\hat{U}_{a1} = \frac{4}{\pi} \cdot U_{dc} \cdot cos \alpha \\
\hat{U}_{av} = \frac{4}{\pi} \cdot U_{dc} \cdot \frac{1}{v} cos(v \cdot \alpha)
\end{gather*}

% ---------------------------------------------------------------------
\subsection{Dreiphasig}
\fig[0.9]{wr_dreiphasig}
\begin{gather*}
u_0(t) = \frac{1}{3} \cdot (u_{a0} + u_{b0} + u_{c0}) \\
u_a(t) = u_{a0}(t) - u_0(t) \\
u_b(t) = u_{b0}(t) - u_0(t) \\
u_c(t) = u_{c0}(t) - u_0(t)
\end{gather*}

\columnbreak
% ---------------------------------------------------------------------
\subsection{Steuerverfahren}
\fig[0.6]{wr_grundfrequenztaktung}
\begin{gather*}
u_{a0}(t) = \frac{4 \cdot U_{dc}}{\pi \cdot 2} \cdot \sum_v^\infty \frac{1}{v} \cdot sin(v \cdot \omega_1 \cdot t) \\
u_a(t) = \frac{4 \cdot U_{dc}}{\pi \cdot 2} \cdot \sum_v^\infty \frac{1}{v} \cdot sin(v \cdot \omega_1 \cdot t) \\
\hat{U}_{a,v=1}(t) = \frac{4 \cdot U_{dc}}{\pi \cdot 2} = 1.273 \cdot \frac{U_{dc}}{2} \\
\hat{U}_{a,v}(t) = \frac{4 \cdot U_{dc}}{\pi \cdot 2} \cdot \frac{1}{v}
\end{gather*}
\begin{achtung}
Ströme der n mal 3 Harmonischen sind 0
\end{achtung}
\begin{gather*}
\hat{I}_{a,v}(t) = \frac{\hat{U}_{a,v}(t)}{\omega \cdot L} \cdot \frac{1}{v} \\
\hat{U}_{a,v=1,rms}(t) = \frac{1}{\sqrt{2}} \frac{4 \cdot U_{dc}}{\pi \cdot 2} = 0.900 \cdot \frac{U_{dc}}{2} \\
\hat{U}_{a,rms}(t) = \sqrt{\frac{8}{9}} \frac{U_{dc}}{2} = 0.943 \cdot \frac{U_{dc}}{2}
\end{gather*}

\columnbreak
\subsubsection{Trägerverfahren dreiphasig}
\fig[0.45]{wr_traegerverfahren}
\begin{gather*}
M_{max} = \frac{2}{\sqrt{3}}
\end{gather*}
Einphasige Brückenschaltung:
\begin{gather*}
\hat{U}_{a,v=1} = M \cdot U_{dc}
\end{gather*}
Mittelpunktschaltung:
\begin{gather*}
\hat{U}_{a,v=1} = M \cdot \frac{U_{dc}}{2}
\end{gather*}
beide mit:
\begin{gather*}
M = \frac{\hat{U}_{sinusf\ddot{o}rmigerSollwert}}{\hat{U}_{Tr\ddot{a}gersignal}}
\end{gather*}
Zu erwartende Frequenzanteile:\\
\begin{symtab}
f & Frequenz Grundschwingung & [Hz] \\
f$_{HL}$ & Schaltfrequenz Halbleiter & [Hz] \\
\end{symtab}
{\setlength{\jot}{1pt}
\begin{gather*}
f_{HL} \pm 2 \cdot f, f_{HL} \pm 4 \cdot f, f_{HL} \pm 6 \cdot f, \\
2 \cdot f_{HL} \pm f, 2 \cdot f_{HL} \pm 3 \cdot f, 2 \cdot f_{HL} \pm 5 \cdot f, \\
3 \cdot f_{HL} \pm 2 \cdot f
\end{gather*}}
Zu erwartende Harmonische:\\
Alle ausser Vielfache von 3.\\
5, 7, 11, 13, 17 \dots

\subsubsection{Drehzeigermodulation}
{\centering\tiny\setlength{\tabcolsep}{1.5pt}\renewcommand{\arraystretch}{1.05}
\begin{tabular}{|c|c|c|c|c|c|c|c|c|c|c|}
\hline
Zustand & $s_a$ & $s_b$ & $s_c$ & $u_{a0}$ & $u_{b0}$ & $u_{c0}$ & $u_a$ & $u_b$ & $u_c$ & $u_0$ \\ \hline
$^0Z$ & -1 & -1 & -1 & $-U_{dc}/2$ & $-U_{dc}/2$ & $-U_{dc}/2$ & 0 & 0 & 0 & $-U_{dc}/2$ \\ \hline
$^1Z$ & 1 & -1 & -1 & $U_{dc}/2$ & $-U_{dc}/2$ & $-U_{dc}/2$ & $2U_{dc}/3$ & $-U_{dc}/3$ & $-U_{dc}/3$ & $-U_{dc}/6$ \\ \hline
$^2Z$ & 1 & 1 & -1 & $U_{dc}/2$ & $U_{dc}/2$ & $-U_{dc}/2$ & $U_{dc}/3$ & $U_{dc}/3$ & $-2U_{dc}/3$ & $U_{dc}/6$ \\ \hline
$^3Z$ & -1 & 1 & -1 & $-U_{dc}/2$ & $U_{dc}/2$ & $-U_{dc}/2$ & $-U_{dc}/3$ & $2U_{dc}/3$ & $-U_{dc}/3$ & $-U_{dc}/6$ \\ \hline
$^4Z$ & -1 & 1 & 1 & $-U_{dc}/2$ & $U_{dc}/2$ & $U_{dc}/2$ & $-2U_{dc}/3$ & $U_{dc}/3$ & $U_{dc}/3$ & $U_{dc}/6$ \\ \hline
$^5Z$ & -1 & -1 & 1 & $-U_{dc}/2$ & $-U_{dc}/2$ & $U_{dc}/2$ & $-U_{dc}/3$ & $-U_{dc}/3$ & $2U_{dc}/3$ & $-U_{dc}/6$ \\ \hline
$^6Z$ & 1 & -1 & 1 & $U_{dc}/2$ & $-U_{dc}/2$ & $U_{dc}/2$ & $U_{dc}/3$ & $-2U_{dc}/3$ & $U_{dc}/3$ & $U_{dc}/6$ \\ \hline
$^7Z$ & 1 & 1 & 1 & $U_{dc}/2$ & $U_{dc}/2$ & $U_{dc}/2$ & 0 & 0 & 0 & $U_{dc}/2$ \\ \hline
\end{tabular}\par}
\begin{gather*}
\vec{u}(t) = u_\alpha + j \cdot u_\beta \\
u_\alpha = \frac{1}{3}(2 \cdot u_a - u_b - u_c) \\
u_\beta = \frac{1}{\sqrt{3}}(u_b - u_c)
\end{gather*}

\columnbreak
{\centering\footnotesize\setlength{\tabcolsep}{4pt}\renewcommand{\arraystretch}{1.1}
\begin{tabular}{|c|c|c|c|c|c|c|c|}
\hline
Zustand & $s_a$ & $s_b$ & $s_c$ & Zeiger & $u_\alpha$ & $u_\beta$ & $|\vec{u}|$ \\ \hline
$^0Z$ & -1 & -1 & -1 & $^0\ul{U}$ & 0 & 0 & 0 \\ \hline
$^1Z$ & 1 & -1 & -1 & $^1\ul{U}$ & $2U_{dc}/3$ & 0 & $2U_{dc}/3$ \\ \hline
$^2Z$ & 1 & 1 & -1 & $^2\ul{U}$ & $U_{dc}/3$ & $U_{dc}/\sqrt{3}$ & $2U_{dc}/3$ \\ \hline
$^3Z$ & -1 & 1 & -1 & $^3\ul{U}$ & $-U_{dc}/3$ & $U_{dc}/\sqrt{3}$ & $2U_{dc}/3$ \\ \hline
$^4Z$ & -1 & 1 & 1 & $^4\ul{U}$ & $-2U_{dc}/3$ & 0 & $2U_{dc}/3$ \\ \hline
$^5Z$ & -1 & -1 & 1 & $^5\ul{U}$ & $-U_{dc}/3$ & $-U_{dc}/\sqrt{3}$ & $2U_{dc}/3$ \\ \hline
$^6Z$ & 1 & -1 & 1 & $^6\ul{U}$ & $U_{dc}/3$ & $-U_{dc}/\sqrt{3}$ & $2U_{dc}/3$ \\ \hline
$^7Z$ & 1 & 1 & 1 & $^7\ul{U}$ & 0 & 0 & 0 \\ \hline
\end{tabular}\par}

\fig[0.6]{wr_raumzeiger_sechseck}

\subsubsection{Optimierte Pulsmuster}
\begin{gather*}
\hat{U}_v = \frac{4 \cdot U_{dc}}{\pi \cdot 2} \cdot \frac{1}{v} \cdot (1 - 2 \cdot cos(v \cdot \alpha_1) + 2 \cdot cos(v \cdot \alpha_2)) \\
M \cdot \frac{4}{\pi} = (1 - 2 \cdot cos(\alpha_1) + 2 \cdot cos(\alpha_2)) % TODO prüfen: laut Skript M * pi/4 = ...
\\
\alpha_2 = arccos\left(\frac{M \cdot \frac{\pi}{4} - 1}{2} + cos(\alpha_1)\right) \\
0 = 1 - 2 \cdot cos(v \cdot \alpha_1) \\
+ 2 \cdot cos\left(v \cdot arccos\left(\frac{M \cdot \frac{\pi}{4} - 1}{2} + cos(\alpha_1)\right)\right)
\end{gather*}
TR harmsol(v,M) für $\alpha_1$ ergebnis kleiner 90 und grösser 0!!

\columnbreak
\fig[0.6]{wr_optimiertes_pulsmuster}

  % !TEX root = ../main.tex
% =====================================================================
\section{Zeiger}
% =====================================================================
\begin{symtab}
$U$ & Spannung & [V] \\
$I$ & Strom & [A] \\
$\hat{I}$ & Stromamplitude & [A] \\
$L_{ph}$ & Phaseninduktivität & [H] \\
$\psi$ & Flussverkettung & [Wb] \\
$\omega$ & Kreisfrequenz & [rad/s] \\
$f$ & Frequenz & [Hz] \\
$t$ & Zeit & [s] \\
$\varphi_i$ & Anfangsphase des Stroms & [rad] \\
$\theta$ & Raumwinkel (Wicklungsachse) & [rad] \\
\end{symtab}

\subsection{Drehfeldzeiger}
\fig[0.6]{drehfeldzeiger}

Drei räumlich um $120^\circ$ versetzte Spulen mit zeitlich um $-120^\circ$
phasenverschobenen Strömen erzeugen ein Drehfeld konstanter Amplitude.\\
Für die drei um $120^\circ$ phasenverschobenen Ströme gilt
\[
i_k(t) = \hat{I}\cos(\omega t + \varphi_k), \qquad \psi_k(t) = L_{ph} i_k(t), \quad k \in \{a,b,c\}.
\]
Die Flussverkettungen der drei Phasen bilden durch geometrische Addition
ein Drehfeld. Die reelle Achse des Koordinatensystems fällt mit der
Wicklungsachse von Phase a zusammen. Die Addition der Momentanwerte liefert
den resultierenden Flussverkettungs-Drehzeiger.

\columnbreak
\fig[0.6]{drehfeldzeiger_fluss}

Mit dem Konventionsfaktor $2/3$ ergibt sich für den Flussverkettungs-Drehzeiger
\[
\vec{\psi} = \frac{2}{3} \sum_{k=a,b,c} \psi_k e^{j\theta_k} = \hat{I} L_{ph} e^{j(\omega t + \varphi_i)}
\]

\subsubsection{Drehfeldzeigertransformation: Dreiphasenebene $\Rightarrow$ Drehzeiger}
Dreiphasige Grössen lassen sich analog zur Flussverkettung als Drehzeiger
in der komplexen Ebene darstellen. Die Transformation gilt allgemein für
Flussverkettungen, Ströme, Spannungen sowie beliebige Dreiphasengrößen $x$.
\begin{gather*}
\vec{\psi}(t) = \frac{2}{3}\Big[\psi_a(t) + \psi_b(t) e^{j120^\circ} + \psi_c(t) e^{-j120^\circ}\Big] \\
\vec{i}(t) = \frac{2}{3}\Big[i_a(t) + i_b(t) e^{j120^\circ} + i_c(t) e^{-j120^\circ}\Big] \\
\vec{u}(t) = \frac{2}{3}\Big[u_a(t) + u_b(t) e^{j120^\circ} + u_c(t) e^{-j120^\circ}\Big] \\
\vec{x}(t) = \frac{2}{3}\Big[x_a(t) + x_b(t) e^{j120^\circ} + x_c(t) e^{-j120^\circ}\Big]
\end{gather*}
Die Drehzeigerdarstellung ermöglicht eine vereinfachte Analyse gegenüber
der Beschreibung mit einzelnen Phasengrössen. Für symmetrische,
sinusförmige Dreiphasensysteme besitzt der Drehzeiger eine konstante Länge
und rotiert mit der Frequenz der Phasengrössen. Seine Spitze beschreibt
dabei einen Kreis in der komplexen Ebene.

\subsubsection{Drehfeldzeigertransformation: Drehzeiger $\Rightarrow$ Zeiger}
Durch Transformation in ein mit $\omega_K$ rotierendes Koordinatensystem wird der
Drehzeiger zu einem zeitlich konstanten Zeiger.
\begin{gather*}
X(t) = \vec{x}(t)\, e^{-j\int \omega_K\, dt}, \\
X(t) = \vec{x}(t)\, e^{-j\omega_K t} \qquad (\omega_K = \text{konstant}), \\
X(t) = X_x + j\,X_y \;(= X_d + j\,X_q).
\end{gather*}
\begin{achtung}
$\omega$ beschreibt die physikalische Drehung\\
$\omega_K$ die Drehung des Bezugssystems.\\
Gleich sind sie \textbf{nur} im synchronen Koordinatensystem
\end{achtung}

\columnbreak
\subsubsection{Drehfeldzeigertransformation: Drehzeiger $\Rightarrow$ Phasengrösse}
Die Rücktransformation ist eindeutig nur bei Berücksichtigung der gleichphasigen Komponente $x_0$.
\vspace{2em}
\begin{align*}
x_a &= x_\alpha + x_0, \\
x_b &= \tfrac{1}{2}(\sqrt{3}x_\beta - x_\alpha) + x_0 \\
x_c &= \tfrac{1}{2}(-\sqrt{3}x_\beta - x_\alpha) + x_0
\end{align*}
\vspace{1em}

  % !TEX root = ../main.tex
% =====================================================================
\section{Halbleiter}
% =====================================================================
\begin{symtab}
E$_{ein}$ & Energie einschaltvorgang & [Ws] \\
E$_{aus}$ & Energie ausschaltvorgang & [Ws] \\
U & Spannung & [V] \\
I & Strom & [A] \\
P$_{vSchalt}$ & Verlustleistung beim Schalten & [W] \\
P$_{vLeit}$ & Verlustleitung beim Leiten & [W] \\
r$_D$ & Widerstand Halbleiter & [$\Omega$] \\
f & Schaltfrequenz & [Hz] \\
T$_{ein}$ & einschaltzeit & [s] \\
T$_{aus}$ & ausschaltzeit & [s] \\
\end{symtab}

\subsection{Graphen}
\vspace*{\fill}
\fig[1.0]{halbleiter_tabelle}

\columnbreak
\fig[0.6]{halbleiter_leistung_frequenz}

% ---------------------------------------------------------------------
\subsection{Schaltverluste}
\fig[0.6]{schaltverluste}
\begin{gather*}
E_{ein} = \int_0^{T_{ein}} u(t) \cdot i(t)\, dt \\
E_{aus} = \int_0^{T_{aus}} u(t) \cdot i(t)\, dt \\
P_{vSchalt} = (E_{ein} + E_{aus}) \cdot f
\end{gather*}

\columnbreak
% ---------------------------------------------------------------------
\subsection{Leitverluste}
\fig[0.6]{leitverluste_ersatz}
\begin{gather*}
P_{vLeit} = \frac{1}{T} \cdot \int_0^T u(t) \cdot i(t)\, dt = U_D \cdot I_{avg} + I_{rms}^2 \cdot r_D
\end{gather*}

% ---------------------------------------------------------------------
\subsection{Kühlung und Halbleitertemperatur}
\begin{symtab}
R$_{th}$ & Thermischer widerstand & [K/W] \\
T & Temperatur & [K, C] \\
P$_{vTot}$ & Verlustleistung Total & [W] \\
\end{symtab}
\fig[0.75]{kuehlung}
\begin{gather*}
\Delta T_{jc} = T_j - T_c = P_{vTot} \cdot R_{thjc} \\
P_{vTot} = P_{vLeit} + P_{vSchalt} \\
\Delta T_{ch} = T_c - T_h = P_{vTot} \cdot R_{thch} \\
\Delta T_{ha} = T_h - T_a = P_{vTot} \cdot R_{thha} \\
\Delta T_{ja} = T_j - T_a = P_{vTot} \cdot R_{thja} \\
R_{thja} = R_{thjc} + R_{thch} + R_{thha} \\
T_j = T_a + \Delta T_{ja}
\end{gather*}

  % !TEX root = ../main.tex
\columnbreak
% =====================================================================
\section{Regelung}
% =====================================================================
\subsection{Regelstruktur Netzseite}
\fig[0.75]{regelstruktur_netzseite}
\fig[0.6]{regelstruktur_entkoppelt}

\subsection{Regelstruktur Uf verfahren}
\fig[0.75]{regelstruktur_uf}

\end{multicols*}

\end{document}
